\documentclass[10pt,a4paper]{amsart}
\usepackage[margin=19mm]{geometry}
\usepackage{amsmath,amssymb,mathtools}
\usepackage[scr=rsfs,cal=euler]{mathalfa}
\usepackage{xfrac}
\usepackage[unicode,hidelinks,bookmarks=false]{hyperref}
\newcommand{\ie}{i.e.\ }
\newcommand{\cf}{cf.\ }
\newcommand{\resp}{respectively\ }
\newcommand{\Subsection}{\subsection}
\providecommand{\nobreakdash}{\nobreak\hbox{-}}
\setlength{\emergencystretch}{2em}
\multlinegap0pt
\usepackage{amssymb}
\usepackage{mathtools}
\usepackage{xcolor}
\usepackage{bm}
\usepackage[shortlabels]{enumitem}
\usepackage{siunitx}
\newtheorem{proposition}{Proposition}
\newcommand{\R}{\mathbb{R}}
\newcommand{\Symm}{\mathbb{S}}
\DeclareMathOperator{\Tr}{Tr}
\title{\LARGE \bf On the Global Optimality of Linear Policies for \\Sinkhorn Distributionally Robust Linear Quadratic Control}
\author{Riccardo Cescon, Andrea Martin, Giancarlo Ferrari-Trecate}
\date{Source: arXiv:2509.00956v2 (CC BY 4.0)}
\begin{document}
\maketitle

\noindent\emph{Source and licence.} \LARGE \bf On the Global Optimality of Linear Policies for \\Sinkhorn Distributionally Robust Linear Quadratic Control. Version arXiv:2509.00956v2 (CC BY 4.0), distributed by arXiv under the Creative Commons Attribution 4.0 International licence, http://creativecommons.org/licenses/by/4.0/, as recorded in the arXiv licence metadata for this exact version and shown on the abstract page https://arxiv.org/abs/2509.00956v2. Full licence notice: \texttt{licenses/arxiv-2509.00956-CC-BY-4.0.txt}.\par\smallskip

\noindent\textit{Editorial scope.} Counted unit: the article's proposition prop:SDP\_reformulation (source0.tex lines 1037--1117). The article's complete original proof of it is delivered with it (source0.tex lines 1118--1247, one \texttt{proof} environment of 130 lines). No mathematics is written, repaired or re-derived: the statement and the proof are the article's own lines, in the article's own notation, and no retained line is substituted or re-flowed.  Retained uncounted as the source-local context the proof needs: lines 557--560.  Nothing was omitted from any retained span, so the package carries no omission record.  No retained line contains a citation, so the wrapper carries no bibliography and no bibliography entry.  No retained line contains a figure environment, an image-inclusion command or drawing code, so no image is delivered and the package declares no asset dependency.  Editorial formatting only, in addition: an AMS \texttt{amsart} preamble that restates the article's own declarations and macros used by the retained lines, namely the environments \texttt{proposition} on separate counters, together with the standard packages those lines need.  The retained environments print in this document as Proposition 1 in order of appearance, whereas the article prints the same items with the numbering of its own full document.  The complete native source of the article as distributed in the arXiv e-print for this exact version is bundled unchanged as \texttt{sources/184208/source0.tex}.
\begin{equation}
\label{eq:trace_lowerbound}
    \underline{d}^\star = \max _{\substack{\mathbf{W}\in\mathcal{G}_W\\ \mathbf{V}\in\mathcal{G}_V}}\ \min _{\substack{\mathbf{U}\in\mathcal{U}^{\text{lin}}\\ \mathbf{q}\in\R^{mT}}}\ \Tr\left(\mathbf{F}_{1} \mathbf{W}+\mathbf{F}_{2} \mathbf{V}\right)+\mathbf{q}^{\top} \mathbf{M} \mathbf{q},
    \end{equation}

\begin{proposition}
\label{prop:SDP_reformulation}
The optimization problem \eqref{eq:trace_lowerbound} is equivalent to the conic program
\begin{align*}
\max \ 
&\Tr(\mathbf{G^{\top} Q G W}) - 
\Tr\big(\mathbf{M}^{-1} F\big)
\\
\text{s.t.} \
&\Gamma \in \mathcal{M}, \ 
F \in \mathbb{S}_{+}^{mT},\mathbf{W} \in \mathbb{S}_{++}^{d(T+1)}, \ 
\mathbf{V} \in \mathbb{S}_{++}^{pT},
\\[1pt]
& E_{x_0} \in \mathbb{S}_{+}^{d}, \ 
E_{w_t} \in \mathbb{S}_{+}^{d}, \ 
E_{v_t} \in \mathbb{S}_{+}^{p}, \ \forall t \in [T],
\\[2pt]
&\begin{aligned}
&\Tr\!\left(\!W_0\! -\! 2E_{x_0}\!+\!\frac{\epsilon}{2}\Sigma^{-1}W_0\!\right)\! -\! \frac{\epsilon}{2}\log\left|E_{x_0}\! - \frac{\epsilon}{4}I\right| \leq \rho_{x_0}
\\
&- \Tr(\hat{X}_0) + \frac{\epsilon}{2}\left(d\log\left(\frac{\epsilon}{2}\right)-\log|\hat{X}_0\Sigma|\right),
\end{aligned}
\\[2pt]
&\begin{aligned}
&\Tr\left(\!W_{t+1}\!\! -\! 2E_{w_t}\!\!+\!\frac{\epsilon}{2}\Sigma^{-1}W_{t+1}\!\right)\!\! -\! \frac{\epsilon}{2}\!\log\!\left|E_{w_t}\!\! -\! \frac{\epsilon}{4}I\right|\! \!\leq\! \rho_{w_t}
\\
&- \Tr(\hat{W}_t) + \frac{\epsilon}{2}\left(d\log\left(\frac{\epsilon}{2}\right)-\log|\hat{W}_t\Sigma|\right),\ \forall t \in [T]
\end{aligned}
\\[2pt]
&\begin{aligned}
&\Tr\left(V_{t}\! -\! 2E_{v_t}\!+\!\frac{\epsilon}{2}\Sigma^{-1}V_{t}\right)\!-\! \frac{\epsilon}{2}\log\left|E_{v_t}\! -\! \frac{\epsilon}{4}I\right| \leq\rho_{v_t}
\\
&- \Tr(\hat{V}_t) + \frac{\epsilon}{2}\left(d\log\left(\frac{\epsilon}{2}\right)-\log|\hat{V}_t\Sigma|\right),\ \forall t \in [T]
\end{aligned}
\\[2pt]
& 
\begin{bmatrix}
\hat{X}_0^{1/2} W_0 \hat{X}_0^{1/2} +\frac{\epsilon^2}{16}I_n & E_{x_0} \\
E_{x_0} & I_n
\end{bmatrix} \succeq 0, 
\\[1pt]
& 
\begin{bmatrix}
\hat{W}_t^{1/2} W_{t+1} \hat{W}_t^{1/2} +\frac{\epsilon^2}{16}I_n & E_{w_t} \\
E_{w_t} & I_n
\end{bmatrix} \succeq 0,\ \forall t \in [T],
\\[1pt]
& 
\begin{bmatrix}
\hat{V}_t^{1/2} V_t \hat{V}_t^{1/2}  +\frac{\epsilon^2}{16}I_p& E_{v_t} \\
E_{v_t} & I_p
\end{bmatrix} \succeq 0, 
\ \forall t \in [T],
\\[1pt]
& 
\setlength\arraycolsep{.9pt}
\begin{bmatrix}
F & \mathbf{H^{\top} Q G W D^{\top}} + \Gamma/2
\\[1pt]
(\mathbf{H^{\top} Q G W D^{\top}} + \Gamma/2)^{\top} & \mathbf{D W D^{\top} + V}
\end{bmatrix}
\!\!\succeq\! 0\,.
% & W_0 \succeq \lambda_{\min}(\hat{X}_0) I, \quad 
% W_{t+1} \succeq \lambda_{\min}(\hat{W}_t) I, \quad 
% V_t \succeq \lambda_{\min}(\hat{V}_t) I, 
% \quad \forall t \in [T-1].
\end{align*}
\noindent
Here, $\mathcal{M}$ denotes the set of all strictly upper block triangular matrices of the form
\[
\Gamma = \begin{bmatrix}
0 & \Gamma_{1,2} & \Gamma_{1,3} & \cdots & \Gamma_{1,T} \\
 & 0 & \Gamma_{2,3} & \cdots & \Gamma_{2,T} \\
 &  & \ddots & \ddots & \vdots \\
&  &  & 0 & \Gamma_{T-1,T} \\
 & &&& 0
\end{bmatrix}
\in \mathbb{R}^{Tm \times Tp}\,,
\]
where $\Gamma_{t,s} \in \mathbb{R}^{m \times p}$ for every $t,s \in \mathbb{Z}$ with $1 \le t < s \le T$.
\end{proposition}

\begin{proof}
The proof relies on dualizing the inner minimization problem in \eqref{eq:trace_lowerbound}. Strong duality holds because the primal problem is feasible and involves only equality constraints, hence any feasible point is in fact a Slater point.
% Note that strong duality holds because the primal problem is trivially feasible and involves only equality constraints, 
% which implies that any feasible point is in fact a Slater point.

In the following we use $\Gamma \in \mathcal{M}$ to denote the Lagrange multiplier of the constraint $\bf U \in \mathcal{U}^{\text{lin}}$, 
which requires all blocks of $\bf U$ above the main diagonal to vanish. 
\noindent
The Lagrangian function of the inner minimization problem in \eqref{eq:trace_lowerbound} is
\begin{equation*}
\mathcal{L}(\mathbf{q}, \mathbf{U}, \Gamma) 
= \Tr\left(\mathbf{F}_{1} \mathbf{W}+\mathbf{F}_{2} \mathbf{V}\right)+\mathbf{q}^{\top} \mathbf{M} \mathbf{q} + \Tr(\mathbf{U}\Gamma^\top)
\end{equation*}
\noindent
Recall now that $\mathbf{R} \succ 0$ and $\mathbf{Q} \succeq 0$, and thus $\mathbf{M} \succ 0$. 
Consequently, $\mathcal{L}$ is minimized by $\mathbf{q}^{\star} = 0$ for any fixed $\bf U$ and $\Gamma$. 
In addition, the partial gradient of $\mathcal{L}$ with respect to $\bf U$ is given by
\[
\frac{\partial \mathcal{L}}{\partial \bf U}
= 2 \mathbf{M U D W D}^{\top} 
 + 2 \mathbf{M U V}
 + 2 \mathbf{H^\top QGWD^\top} + \Gamma\,.
\]
\noindent
Recall also that $\mathbf{V} \in \mathcal{G}_{V}$ is strictly positive, which implies that $\mathbf{D W D^{\top} }+ \mathbf{V} \succ 0$ is invertible. 
As we already know that $\mathbf{M}\succ 0$ is invertible as well, 
$\mathcal{L}$ is minimized by
\[
\mathbf{U}^{\star} 
= -\mathbf{M}^{-1} 
   (\mathbf{H^{\top} Q G W D^{\top}} + \Gamma/2)
   (\mathbf{D W D^{\top} + V})^{-1}
\]
for any fixed $\Gamma$.

\noindent
Substituting both $\bf q^{\star}$ and $\bf U^{\star}$ into $\mathcal{L}$ yields the dual objective function
\[
\begin{aligned}
&g(\Gamma)
= \mathcal{L}(\mathbf{q^{\star}, U}^{\star}, \Gamma)
= \Tr(\mathbf{G^{\top} Q G W})+
\\
&
\begin{multlined}
- \Tr\Big(
    (\mathbf{M}^{-1}
    (\mathbf{H^{\top} Q G W D^{\top}} + \Gamma/2)
    (\mathbf{D W D^{\top} + V})^{-1}
    \\
    (\mathbf{H^{\top} Q G W D^{\top}} + \Gamma/2)^{\top}
   \Big)\,.
\end{multlined}
\end{aligned}
\]

\noindent
The dual of the inner minimization problem in \eqref{eq:trace_lowerbound} is thus given by 
$\max_{\Gamma \in \mathcal{M}} g(\Gamma)$.
To linearize the dual function $g(\Gamma)$, we introduce the auxiliary variable $F \in \mathbb{S}_{+}^{mT}$ 
subject to the matrix inequality
\[
\begin{multlined}
F \succeq 
  (\mathbf{H^{\top} Q G W D^{\top}} + \Gamma/2)
    (\mathbf{D W D^{\top} + V})^{-1}
    \\
    (\mathbf{H^{\top} Q G W D^{\top}} + \Gamma/2)^{\top}\,.
\end{multlined}
\]
We can reformulate the previous inequality as an LMI using the Schur complement. Then, the dual problem rewrites as
\begin{equation}
\label{eq:inner_dual}
\begin{aligned}
\max \ \ &\Tr(\mathbf{G^{\top} Q G W}) - 
\Tr\big(\mathbf{M}^{-1} F\big)
\\[1pt]
\text{s.t.} \ \ & 
\Gamma \in \mathcal{M}, \ 
F \in \mathbb{S}_{+}^{mT},
\\[1pt]
& 
\!\!\setlength\arraycolsep{.9pt}
\begin{bmatrix}
F & \mathbf{H^{\top} Q G W D^{\top}} + \Gamma/2
\\[1pt]
(\mathbf{H^{\top} Q G W D^{\top}} + \Gamma/2)^{\top} & \mathbf{D W D^{\top} + V}
\end{bmatrix}
\!\!\succeq\!\!0\,.
\end{aligned}
\end{equation}
Substituting the strong dual \eqref{eq:inner_dual} in the inner minimization in \eqref{eq:trace_lowerbound}, we obtain
\begin{align*}
  \max \ \ &\Tr(\mathbf{G^{\top} Q G W}) - 
\Tr\big(\mathbf{M}^{-1} F\big)
\\[1pt]
\text{s.t.} \ \ & 
\Gamma \in \mathcal{M}, \ 
F \in \mathbb{S}_{+}^{mT},\mathbf{W} \in \mathbb{S}_{++}^{d(T+1)}, \ 
\mathbf{V} \in \mathbb{S}_{++}^{pT},
\\[1pt]
& 
\!\!\setlength\arraycolsep{.9pt}
\begin{bmatrix}
F & \mathbf{H^{\top} Q G W D^{\top}} + \Gamma/2
\\[1pt]
(\mathbf{H^{\top} Q G W D^{\top}} + \Gamma/2)^{\top} & \mathbf{D W D^{\top} + V}
\end{bmatrix}
\!\!\succeq\! 0,
\\[1pt]
&G_\epsilon(\hat{X}_0, W_0) \leq \rho_{x_0},
\\
&G_\epsilon(\hat{W}_{t+1}, W_{t+1}) \leq \rho_{w_t},\ G_\epsilon(\hat{V}_t, V_t) \leq \rho_{v_t} \ \forall t \in [T]\,.
\end{align*}
We conclude the proof by reformulating the entropy-regularized Gelbrich constraints. For simplicity, consider the constraint $G_\epsilon(\hat{X}_0, W_0) \leq \rho_{x_0}$; the reasoning applies \textit{mutatis mutandis} to the others. Consider an auxiliary variable $E_{x_0}\in\Symm^n_{+}$ subject to the matrix inequality $E_{x_0}^2 \preceq \hat{X}_0^{1/2} W_0 \hat{X}_0^{1/2} +\frac{\epsilon^2}{16}I_n$. This inequality can be recast as $E_{x_0} \preceq (\hat{X}_0^{1/2} W_0 \hat{X}_0^{1/2} +\frac{\epsilon^2}{16}I_n)^\frac{1}{2}$. With this auxiliary variable, we can equivalently express the nonlinear inequality constraint as
\begin{equation*}
\begin{multlined}
   \Tr\left(W_0 - 2E_{x_0}+\frac{\epsilon}{2}\Sigma^{-1}W_0\right) - \frac{\epsilon}{2}\log\left|E_{x_0} - \frac{\epsilon}{4}I\right| \leq \rho_{x_0}
   \\
   - \Tr(\hat{X}_0) + \frac{\epsilon}{2}\left(d\log\left(\frac{\epsilon}{2}\right)-\log|\hat{X}_0\Sigma|\right)\,.
\end{multlined}
\end{equation*}
If the constraint is satisfied with $E_{x_0}$ it will be also with the original variable since the trace and the log-determinant are monotone operators. Finally, with a Schur complement argument the inequality $E_{x_0}^2 \preceq \hat{X}_0^{1/2} W_0 \hat{X}_0^{1/2} +\frac{\epsilon^2}{16}I_n$ can be equivalently rewritten as
\begin{equation*}
    \begin{bmatrix}
\hat{X}_0^{1/2} W_0 \hat{X}_0^{1/2} +\frac{\epsilon^2}{16}I_n & E_{x_0} \\
E_{x_0} & I_n
\end{bmatrix} \succeq 0\,.
\end{equation*}
\end{proof}

\end{document}
